% ============================================================
%  Hybrid BRD v3.0 — Decision-Making Needs Expression
%  Project: AI-Based IDS/IPS — Port Scan & Reconnaissance
%  Includes: ML Detection + Honeypot (first-class) + MTD (response action)
%  Module: Artificial Intelligence — Semester 2
%  Version: 3.0 — April 2026
% ============================================================
\documentclass[12pt,a4paper]{article}

% --- Packages ---
\usepackage[utf8]{inputenc}
\usepackage[T1]{fontenc}
\usepackage[english]{babel}
\usepackage{geometry}
\usepackage{enumitem}
\usepackage{booktabs}
\usepackage{array}
\usepackage{xcolor}
\usepackage{titlesec}
\usepackage{fancyhdr}
\usepackage{hyperref}
\usepackage{tabularx}
\usepackage{multirow}
\usepackage{graphicx}

\geometry{margin=2.5cm}

% --- Colors ---
\definecolor{sectionblue}{RGB}{15, 23, 42}
\definecolor{accentcyan}{RGB}{6, 182, 212}
\definecolor{accentgreen}{RGB}{16, 185, 129}
\definecolor{warnorange}{RGB}{234, 179, 8}
\definecolor{lightgray}{RGB}{241, 245, 249}

% --- Section styling ---
\titleformat{\section}
  {\large\bfseries\color{sectionblue}}
  {\thesection.}{0.5em}{}[\vspace{-0.5em}\textcolor{accentcyan}{\rule{\textwidth}{0.4pt}}]

\titleformat{\subsection}
  {\normalsize\bfseries\color{sectionblue}}
  {\thesubsection.}{0.5em}{}

% --- Header/Footer ---
\pagestyle{fancy}
\fancyhf{}
\fancyhead[L]{\small\textcolor{gray}{Hybrid BRD v3.0 --- IDS/IPS Port Scan \& Recon}}
\fancyhead[R]{\small\textcolor{gray}{AI Module --- S2}}
\fancyfoot[C]{\small\thepage}
\renewcommand{\headrulewidth}{0.4pt}

% --- Table column types ---
\newcolumntype{L}[1]{>{\raggedright\arraybackslash}p{#1}}
\newcolumntype{C}[1]{>{\centering\arraybackslash}p{#1}}

% --- Highlight box ---
\newcommand{\infobox}[2]{%
  \medskip
  \noindent
  \fcolorbox{accentcyan}{lightgray}{%
    \parbox{\dimexpr\textwidth-2\fboxsep-2\fboxrule}{%
      \textbf{\textcolor{accentcyan}{#1}}\\[0.3em]
      #2
    }%
  }%
  \medskip
}

\begin{document}

% ============================================================
%  TITLE BLOCK
% ============================================================
\begin{center}
  {\Large\bfseries Hybrid BRD v3.0 --- Decision-Making Needs Expression}\\[0.6em]
  {\large AI-Based Intrusion Detection \& Response System}\\[0.3em]
  {\large Detection of \textit{Port Scan \& Network Reconnaissance} Attacks}\\[0.3em]
  {\normalsize Integrating ML Detection, Honeypot Deception \& Moving Target Defense}\\[1.2em]
  \begin{tabular}{rl}
    \textbf{Module:}   & Artificial Intelligence --- Semester 2 \\
    \textbf{Team:}     & 5 members (Engineering --- ENSA) \\
    \textbf{Version:}  & Hybrid BRD v3.0 \\
    \textbf{Date:}     & April 2026 \\
    \textbf{Severity:} & MEDIUM-HIGH \\
  \end{tabular}
\end{center}

\vspace{0.5em}

% --- Project Team ---
\begin{table}[h!]
\centering
\small
\begin{tabular}{C{3cm} C{3cm} C{3cm} C{3cm} C{3cm}}
  \toprule
  \textbf{El Kartouti Anas} & \textbf{Lekhbioui Adil} & \textbf{Moulay Anas} & \textbf{Hammoubel El Yazid} & \textbf{Nafia Khadija} \\
  Member & Member & Member & Member & Member \\
  \bottomrule
\end{tabular}
\end{table}

\vspace{1em}

% ============================================================
%  1. EXECUTIVE SUMMARY
% ============================================================
\section{Executive Summary}

This document presents the formal decision-making needs expression for the
design, implementation, testing, and validation of an Artificial Intelligence
model dedicated to the \textbf{real-time detection and automated mitigation} of
Port Scanning and Network Reconnaissance attacks within a controlled LAN
environment.

Reconnaissance constitutes the first phase of virtually every cyberattack kill
chain: an adversary probes the network to map active hosts, enumerate open
ports, and fingerprint running services before launching a targeted intrusion.
Early detection and containment at this phase prevents the full materialization
of subsequent, more destructive attack stages.

This Hybrid BRD v3.0 extends the project scope with three complementary
defense layers:

\begin{enumerate}[nosep,leftmargin=1.5em]
  \item \textbf{AI/ML Detection Layer:} A supervised classification pipeline
        (Random Forest + XGBoost comparison + Isolation Forest for slow scans)
        trained on labeled IDS benchmark datasets.
  \item \textbf{Honeypot Deception Layer:} Decoy services deployed via
        Cowrie/Dionaea to attract, trap, and profile scanning adversaries ---
        with honeypot interaction data fed as a real-time feature to the ML model.
  \item \textbf{Moving Target Defense (MTD) Response Layer:} Upon confirmed
        scan detection, the system triggers dynamic port reassignment to
        invalidate the attacker's reconnaissance data.
\end{enumerate}

Together, these three pillars --- detection, deception, and surface mutation ---
form a layered, intelligent security posture.

% ============================================================
%  2. PROBLEM STATEMENT
% ============================================================
\section{Problem Statement}

Port scanning in its various forms (SYN scan, ACK scan, UDP scan, XMAS scan,
FIN scan) produces traffic patterns that differ fundamentally from legitimate
host-to-host communication. An adversary conducting a scan sends probes to
hundreds or thousands of ports in succession, producing an anomalous
distribution of half-open connections, RST responses, and unanswered packets.

Despite being a well-understood technique, \textbf{detecting stealthy or
slow-rate scanning} --- where the adversary deliberately spaces probes to
evade threshold-based detection --- remains a significant challenge. Such
low-and-slow patterns require behavioral analysis over extended time windows
rather than instantaneous packet-level inspection.

A further critical dimension of the problem is the \textbf{static, predictable
nature of conventional network configurations}. Fixed IP addresses, stable
port assignments, and consistent service layouts give adversaries a reliable
map to exploit. Once a scan is complete, the attacker possesses a stable,
actionable picture of the infrastructure --- a problem that no detection system
alone can fully address.

Additionally, once scanning activity is detected, organizations currently
lack systematic means to \textbf{observe and profile} the adversary's intent
and methodology in a controlled, risk-free manner.

The problem statement can thus be formulated as follows:

\begin{quote}
\itshape
``How can we design, train, and evaluate a Machine Learning model capable of
classifying network traffic with high reliability to detect
\textbf{Port Scan \& Reconnaissance} attacks --- including slow and stealthy
variants --- while integrating honeypot-based attacker profiling and
post-detection surface mutation to minimize the value of collected
reconnaissance data?''
\end{quote}


% ============================================================
%  3. PROJECT OBJECTIVES
% ============================================================
\section{Project Objectives}

\begin{table}[h!]
\centering
\small
\begin{tabular}{L{4cm} L{11cm}}
  \toprule
  \textbf{Objective} & \textbf{Description} \\
  \midrule
  Primary Detection     & Train and deploy an AI model capable of detecting port scanning and network reconnaissance activity in real time. \\
  Scan Type Coverage    & Detect and classify SYN Scan (stealth), TCP Connect, UDP Scan, ACK Scan, XMAS Scan, and slow/distributed variants. \\
  Honeypot Deception    & Deploy decoy services (Cowrie/Dionaea) that mimic legitimate targets, trap scanning adversaries, and enrich attacker profiling through logged interactions. \\
  Moving Target Defense & Upon confirmed scan detection, trigger dynamic port reassignment to invalidate the attacker's reconnaissance map. \\
  Attacker Profiling    & Build a structured profile per detected threat actor: network identity, scan methodology, targeted scope, behavioral timeline, and honeypot interaction history. \\
  Real-time Alerting    & Trigger graduated alerts classified by scan type, confidence score, and severity level. \\
  Automated Response    & Enforce automated protective measures --- including network-level blocking and MTD port rotation --- upon confirmed threat detection. \\
  \bottomrule
\end{tabular}
\end{table}


% ============================================================
%  4. FUNCTIONAL REQUIREMENTS
% ============================================================
\section{Functional Requirements}

\subsection{System Inputs}

\begin{itemize}[nosep,leftmargin=1.5em]
  \item Raw TCP/UDP packets captured on the network interface, filtered for
        partial or incomplete handshake patterns indicative of probing behavior.
  \item Labeled training datasets: CICIDS2017 PortScan subset ($\sim$158,000
        flows) and UNSW-NB15 Reconnaissance category ($\sim$13,000 records).
  \item Sliding time-window aggregated flow features computed per source IP
        over configurable observation windows (10\,s and 60\,s).
  \item Network metadata: source/destination IP, MAC addresses, ports,
        protocol, TCP flag combinations, TTL values, and timestamps.
  \item Honeypot interaction logs: connection attempts, payloads, source
        identifiers, and session metadata from decoy nodes.
\end{itemize}

\subsection{System Outputs}

\begin{itemize}[nosep,leftmargin=1.5em]
  \item Real-time dashboard: per-source port probe counts, scan heat map,
        active threat list, honeypot interaction events.
  \item Per-time-window classification: \textsc{Benign} or
        \textsc{Port\_Scan}, with inferred scan type and confidence score.
  \item Attacker profile card: source IP and MAC, scan type, ports and hosts
        targeted, first/last seen timestamps, honeypot interaction flag,
        threat level.
  \item Graduated alert notification with scan type label, severity score,
        and contextual metadata.
  \item Automated network-level block rule applied to the confirmed
        attacker's source address, with the event logged to the incident
        database.
  \item MTD rotation trigger: upon detection, the system initiates a port
        reassignment cycle to invalidate the attacker's reconnaissance data.
  \item \textit{[Optional]} Honeypot redirect: scanner traffic rerouted to a
        decoy listener via \texttt{iptables REDIRECT}.
\end{itemize}


% ============================================================
%  5. DATASET SPECIFICATION
% ============================================================
\section{Dataset Specification}

\begin{table}[h!]
\centering
\small
\begin{tabular}{L{3.2cm} L{4cm} C{3cm} L{3cm} C{2.5cm}}
  \toprule
  \textbf{Dataset} & \textbf{Relevant Subset} & \textbf{Samples (approx.)} & \textbf{Source} & \textbf{Role} \\
  \midrule
  CICIDS2017       & PortScan category         & $\sim$158,000 flows   & UNB / CIC (Canada)    & PRIMARY \\
  UNSW-NB15        & Reconnaissance category   & $\sim$13,000 records  & UNSW (Australia)      & SECONDARY \\
  \bottomrule
\end{tabular}
\end{table}

\infobox{Dataset Rationale}{%
  \textbf{CICIDS2017:} Modern IDS benchmark with labeled attack flows and
  realistic traffic patterns. Widely cited in current IDS research.\\
  \textbf{UNSW-NB15:} Adds diversity --- different capture environment and
  attack representation --- to improve model generalization.\\
  \textbf{KDD-99 (1999):} Explicitly \textbf{REJECTED} --- universally
  considered obsolete in IDS literature since 2010.\\
  \textbf{Live Nmap simulations} will be used for \textit{system validation
  only} (not model training) --- see Section~\ref{sec:timeline}, Phase 3.
}


% ============================================================
%  6. KEY NETWORK FEATURES
% ============================================================
\section{Key Network Features for the ML Model}

\begin{table}[h!]
\centering
\small
\begin{tabular}{L{4cm} L{5.5cm} L{4.5cm}}
  \toprule
  \textbf{Feature Name} & \textbf{Description} & \textbf{Relevance} \\
  \midrule
  Distinct Dst Ports/IP  & \# unique destination ports per source          & Critical --- core scanning signature \\
  SYN Flag Count         & SYN sent without completing handshake           & Critical --- stealth SYN scan \\
  RST Flag Count         & RST responses from closed ports                 & Very High --- port probe rejection \\
  Flow Duration          & Duration of each individual probe connection    & High --- very short in fast scans \\
  Total Fwd Packets      & Packets per flow from scanner                   & High --- typically 1 per probe \\
  IAT Mean               & Mean inter-arrival time between probes          & High --- slow scans have large IAT \\
  Port Range Entropy     & Shannon entropy of destination ports            & High --- randomized vs sequential scan \\
  ACK Flag Count         & ACK packets without prior SYN                   & Medium --- ACK scan detection \\
  Unique Dst IPs/Src     & \# distinct destination hosts contacted         & Medium --- host discovery phase \\
  Bwd Packet Length      & Response packet size (RST vs SYN-ACK ratio)     & Medium --- bidirectional profile \\
  TTL Value              & Time-to-live in IP header                       & Low --- OS fingerprinting attempts \\
  \midrule
  \textbf{Honeypot Interaction Flag} \newline\textit{[NEW v3.0]}
                         & Binary: did source IP contact a decoy node?     & \textbf{Critical} --- high-confidence threat signal \\
  \bottomrule
\end{tabular}
\end{table}

\infobox{Feature Engineering Notes}{%
  \textbf{Port Range Entropy:} Shannon entropy of destination ports ---
  high entropy signals randomized scans. RF-friendly, peer-cited.\\
  \textbf{Honeypot Interaction Flag} \textit{[NEW v3.0]:} Binary indicator
  set to 1 when the source IP has contacted a honeypot decoy. No legitimate
  user ever contacts a honeypot, making this a near-certain threat
  confirmation signal.\\
  Compute all features per \texttt{(src\_ip, 10s window)} AND
  \texttt{(src\_ip, 60s window)} independently.\\
  SMOTE applied on training fold only --- never on validation or test
  splits --- to prevent data leakage.
}


% ============================================================
%  7. MACHINE LEARNING MODEL STRATEGY
% ============================================================
\section{Machine Learning Model Strategy}

\begin{table}[h!]
\centering
\small
\begin{tabular}{L{4.2cm} L{10.5cm}}
  \toprule
  \textbf{Component} & \textbf{Detail} \\
  \midrule
  Primary Model         & Random Forest Classifier --- high-dimensional tabular data,
                          natural feature importance ranking, robust to class imbalance. \\
  Comparison Model      & XGBoost --- gradient boosting compared on time-windowed
                          aggregated features as a high-performance ensemble benchmark. \\
  Slow Scan Layer       & Isolation Forest (unsupervised) on 60-second sliding window
                          per source IP to capture low-and-slow probing patterns. \\
  Feature Engineering   & Time-window aggregation: distinct ports/IP/10s, SYN-ratio,
                          IAT statistics, Port Range Entropy, Honeypot Interaction Flag
                          over rolling windows. \\
  Imbalance Handling    & SMOTE oversampling on the \textsc{Port\_Scan} minority class
                          in the training set only. \\
  Hyperparameter Tuning & GridSearchCV / RandomizedSearchCV on RF and XGBoost ---
                          quantify performance gains vs default parameters. \\
  Evaluation Metrics    & Accuracy, Precision, Recall, F1-Score, AUC-ROC, Detection
                          Rate per scan type, Honeypot True Positive Rate. \\
  Validation Strategy   & Stratified K-Fold ($k=10$) + temporal holdout (last 20\% by
                          chronological order) --- prevents data leakage. \\
  v1 Output             & Binary: \textsc{Benign} vs \textsc{Port\_Scan} with
                          confidence score. \\
  v2+ Output            & Multi-class: SYN / UDP / ACK / XMAS / Connect scan type. \\
  \bottomrule
\end{tabular}
\end{table}


% ============================================================
%  8. MTD & HONEYPOT COMPONENT SPECIFICATIONS
% ============================================================
\section{MTD \& Honeypot Component Specifications}

\subsection{Moving Target Defense (MTD) --- Response Action}

The MTD component is activated as a \textbf{post-detection response action}.
When the AI engine confirms a port scan with high confidence ($\geq 0.85$),
the MTD engine triggers a port rotation cycle to invalidate the attacker's
reconnaissance map.

\begin{table}[h!]
\centering
\small
\begin{tabular}{L{3.5cm} L{11cm}}
  \toprule
  \textbf{Component} & \textbf{Description} \\
  \midrule
  Mutation Scope    & Dynamic rotation of service port assignments on protected
                      hosts within the network. \\
  Rotation Trigger  & Event-based: triggered automatically upon confirmed scan
                      detection by the AI engine ($\text{confidence} \geq 0.85$). \\
  Coordination      & All legitimate nodes are synchronized with the current port
                      assignment table via secure configuration push. \\
  Scope Limitation  & MTD operates as a \textbf{response mechanism only} --- it does
                      not feed data into the ML model. AI-integrated MTD (using MTD
                      state as an ML feature) is deferred to future work. \\
  \bottomrule
\end{tabular}
\end{table}

\infobox{MTD Scope Decision}{%
  MTD is implemented as a \textbf{simple response action} (rotate ports after
  detection), not as an AI-integrated feature. The ``MTD Port Delta'' feature
  proposed in earlier versions has been \textbf{deferred to the Perspectives /
  Future Work chapter} of the final report. Rationale: existing datasets
  (CICIDS2017, UNSW-NB15) were captured on static networks and contain no
  MTD-aware traffic patterns, making model training for this feature
  infeasible within the project timeline.
}

\subsection{Honeypot Integration --- First-Class Component}

Honeypot nodes are deployed within the network topology as high-fidelity decoy
systems. They present a realistic appearance to any scanning adversary while
containing no operational data or functionality. All traffic directed to a
honeypot is inherently unauthorized and constitutes an unambiguous threat signal.

\begin{table}[h!]
\centering
\small
\begin{tabular}{L{3.5cm} L{11cm}}
  \toprule
  \textbf{Component} & \textbf{Description} \\
  \midrule
  Honeypot Type       & Low-to-medium interaction honeypots emulating common services
                        (SSH, HTTP, FTP, Telnet) to attract and engage scanning activity. \\
  Framework           & \textbf{Cowrie} (SSH/Telnet) or \textbf{Dionaea} (multi-service),
                        or custom lightweight Python socket listeners. \\
  Placement Strategy  & Distributed across the network topology to maximize the
                        probability of contact during any broad scanning campaign. \\
  AI Contribution     & Honeypot interaction events --- connection attempts, payloads,
                        source IPs --- are ingested as the \textbf{Honeypot Interaction
                        Flag} feature, significantly boosting detection confidence
                        for confirmed threat actors. \\
  Profiling Output    & Each honeypot interaction is logged and associated with the
                        corresponding attacker profile, enriching the threat intelligence
                        record with behavioral indicators, service targeting patterns,
                        and session-level metadata. \\
  \bottomrule
\end{tabular}
\end{table}

\subsection{Implementation Effort Estimate}

\begin{table}[h!]
\centering
\small
\begin{tabular}{L{5cm} L{5cm} C{3cm}}
  \toprule
  \textbf{Component} & \textbf{Tool} & \textbf{Effort} \\
  \midrule
  Honeypot deployment (Cowrie)     & Docker: \texttt{cowrie/cowrie}         & $\sim$1 hour \\
  Banner/service configuration     & Edit \texttt{cowrie.cfg}               & $\sim$30 min \\
  Log parser → feature pipeline    & Python script (JSON → Pandas)          & $\sim$2 hours \\
  Dashboard integration            & Add Honeypot Activity tab              & $\sim$2 hours \\
  MTD rotation engine              & Python + \texttt{iptables}             & $\sim$3 hours \\
  MTD state synchronization        & Config push to legitimate nodes        & $\sim$2 hours \\
  \bottomrule
\end{tabular}
\end{table}


% ============================================================
%  9. NON-FUNCTIONAL REQUIREMENTS
% ============================================================
\section{Non-Functional Requirements}

\begin{table}[h!]
\centering
\small
\begin{tabular}{L{4cm} L{6.5cm} C{3.5cm}}
  \toprule
  \textbf{Requirement} & \textbf{Criterion} & \textbf{Target Value} \\
  \midrule
  Fast Scan Detection      & Detection time for aggressive Nmap scan         & $< 8$ seconds \\
  Slow Scan Detection      & Detection time for stealthy / slow scan         & $< 90$\,s (sliding window) \\
  Model F1-Score           & On CICIDS2017 PortScan test subset              & $\geq 88\%$ \\
  Model Accuracy           & Overall classification on test set              & $\geq 90\%$ \\
  False Positive Rate      & Normal connections flagged as scans             & $< 6\%$ \\
  Dashboard Refresh        & Scanner activity \& honeypot heat map update    & $\leq 5$ seconds \\
  Blocking Response        & Time from detection to \texttt{iptables} DROP   & $< 15$ seconds \\
  MTD Rotation Latency     & Time from detection trigger to completed rotation & $< 30$ seconds \\
  Honeypot Detection Rate  & Proportion of scanning IPs that interact with
                             at least one honeypot                           & $\geq 60\%$ \\
  \bottomrule
\end{tabular}
\end{table}


% ============================================================
%  10. SYSTEM ARCHITECTURE VISION
% ============================================================
\section{System Architecture Vision}

\subsection{Capture Layer}
Scapy BPF filter on LAN interface (SYN-only or full TCP/UDP).
Rolling aggregation per source IP over 10\,s and 60\,s windows.

\subsection{Feature Layer}
Pandas windowed aggregation $\rightarrow$ 12-feature vector per
\texttt{(src\_ip, time\_window)}. Port Range Entropy, Honeypot Interaction
Flag. SMOTE applied on training split only.

\subsection{Detection Layer}
\begin{itemize}[nosep,leftmargin=1.5em]
  \item \textbf{v1:} RF binary classifier $\rightarrow$ \textsc{Benign} /
        \textsc{Port\_Scan} + confidence score.
  \item \textbf{v2+:} Multi-class scan type (SYN / UDP / ACK / XMAS / Connect).
  \item \textbf{Parallel:} Isolation Forest on 60\,s window for slow scan
        anomaly detection.
\end{itemize}

\subsection{Response Layer}
\begin{itemize}[nosep,leftmargin=1.5em]
  \item \textbf{Alert:} Severity score + scan type label $\rightarrow$ event store.
  \item \textbf{Blocking:} \texttt{iptables} DROP via Python subprocess
        $\rightarrow$ audit log.
  \item \textbf{MTD:} Port rotation triggered upon confirmed detection
        ($\text{confidence} \geq 0.85$).
  \item \textbf{Honeypot redirect} \textit{[optional]:} Scanner traffic
        rerouted to decoy listener via \texttt{iptables REDIRECT}.
\end{itemize}

\subsection{Observability Layer}
Real-time dashboard (Flask + Socket.IO + D3.js):
\begin{itemize}[nosep,leftmargin=1.5em]
  \item Active scanners list with per-IP threat level.
  \item Port activity heat map.
  \item Attacker profile (IP, MAC, ports probed, scan type, first/last seen).
  \item Alert feed with confidence, severity, and scan type label.
  \item Honeypot interaction event log.
\end{itemize}


% ============================================================
%  11. TEAM ROLE ASSIGNMENTS
% ============================================================
\section{5-Person Team Role Assignments}

Each role maps to a deliverable that contributes to the AI lifecycle
documentation grade. Every member owns a graded chapter of the final report.

\begin{table}[h!]
\centering
\small
\begin{tabular}{C{1cm} L{3.5cm} L{3cm} L{6.5cm}}
  \toprule
  \textbf{Role} & \textbf{Title} & \textbf{Assigned To} & \textbf{Key Deliverables} \\
  \midrule
  R1 & ML Pipeline Lead & Lekhbioui Adil &
       Dataset preprocessing, SMOTE, RF + XGBoost training,
       GridSearchCV tuning, comparison table.
       \newline\textit{Report: Data Preparation + Modeling.} \\
  \midrule
  R2 & Anomaly \& Slow Scan Specialist & TBC &
       Isolation Forest on 60\,s window, contamination tuning,
       integration with RF confidence in unified alert score.
       \newline\textit{Report: Unsupervised layer + evaluation.} \\
  \midrule
  R3 & Capture \& Feature Pipeline & TBC &
       Scapy listener with BPF filter, rolling aggregation
       windows (10\,s, 60\,s), real-time feature vector output.
       \newline\textit{Report: Data Understanding + architecture.} \\
  \midrule
  R4 & Dashboard, Honeypot \& MTD & TBC &
       Flask + Socket.IO + D3.js dashboard, Cowrie deployment,
       honeypot log parser, MTD rotation engine.
       \newline\textit{Report: Deployment + demo recording.} \\
  \midrule
  R5 & Attack Simulation \& Docs & TBC &
       Nmap scan types against test VM, per-scan detection table,
       literature review, ethical considerations, French abstract.
       \newline\textit{Report: Problem Statement + Evaluation.} \\
  \bottomrule
\end{tabular}
\end{table}


% ============================================================
%  12. WEEK-BY-WEEK TIMELINE
% ============================================================
\section{Week-by-Week Execution Timeline}
\label{sec:timeline}

\subsection*{Phase 1 --- Design \& Setup (Week 1)}

\begin{table}[h!]
\centering
\small
\begin{tabular}{C{1.2cm} L{7.5cm} C{1.5cm} L{3.5cm}}
  \toprule
  \textbf{Week} & \textbf{Activities} & \textbf{Owner} & \textbf{Deliverable} \\
  \midrule
  Wk 1 & VM creation (Ubuntu IDS + honeypot, Kali attacker), network isolation, Git repo, role confirmation & R3+R5 & Network diagram, Git repo \\
  Wk 1 & Download CICIDS2017 + UNSW-NB15; EDA: class distribution, missing values, descriptive statistics & R1 & Data understanding report \\
  Wk 1 & Deploy Cowrie honeypot on dedicated VM, configure fake banners & R4 & Honeypot node online \\
  \bottomrule
\end{tabular}
\end{table}

\subsection*{Phase 2 --- Implementation (Weeks 2--4)}

\begin{table}[h!]
\centering
\small
\begin{tabular}{C{1.2cm} L{7.5cm} C{1.5cm} L{3.5cm}}
  \toprule
  \textbf{Week} & \textbf{Activities} & \textbf{Owner} & \textbf{Deliverable} \\
  \midrule
  Wk 2 & Data preprocessing: SMOTE, time-window aggregation (10\,s/60\,s), Port Range Entropy & R1+R3 & Clean dataset + feature set \\
  Wk 2 & Scapy live capture: BPF filter, rolling aggregation, real-time feature vector output & R3 & \texttt{capture\_pipeline.py} \\
  Wk 3 & Train RF and XGBoost: baseline $\rightarrow$ GridSearchCV tuning $\rightarrow$ model comparison table & R1 & Trained \texttt{.pkl} + comparison table \\
  Wk 3 & Isolation Forest: contamination tuning, integration with RF confidence scores & R2 & \texttt{slow\_scan\_detector.py} \\
  Wk 3 & Honeypot log parser: extract interaction events, generate Honeypot Interaction Flag & R4 & \texttt{honeypot\_parser.py} \\
  Wk 4 & Flask/D3.js dashboard: heat map, active scanners, alert feed, attacker profile, honeypot tab & R4 & \texttt{dashboard/} module \\
  Wk 4 & \texttt{iptables} blocking engine + MTD rotation engine (port reassignment on detection) & R3+R4 & \texttt{response\_engine.py} \\
  \bottomrule
\end{tabular}
\end{table}

\subsection*{Phase 3 --- Testing (Week 5)}

\begin{table}[h!]
\centering
\small
\begin{tabular}{C{1.2cm} L{7.5cm} C{1.5cm} L{3.5cm}}
  \toprule
  \textbf{Week} & \textbf{Activities} & \textbf{Owner} & \textbf{Deliverable} \\
  \midrule
  Wk 5 & Attack simulation: Nmap SYN/Connect/UDP/XMAS/ACK scans against test VM & R5 & Per-scan detection rate table \\
  Wk 5 & False positive testing: legitimate traffic scenarios (file transfers, browsing) & R2+R5 & FPR measurement report \\
  Wk 5 & Validate honeypot engagement: measure \% of scanners that hit decoy nodes & R4+R5 & Honeypot engagement log \\
  Wk 5 & Validate MTD rotation: confirm port reassignment completes within 30\,s & R3+R4 & MTD effectiveness report \\
  \bottomrule
\end{tabular}
\end{table}

\subsection*{Phase 4 --- Validation \& Report (Weeks 6--7)}

\begin{table}[h!]
\centering
\small
\begin{tabular}{C{1.2cm} L{7.5cm} C{1.5cm} L{3.5cm}}
  \toprule
  \textbf{Week} & \textbf{Activities} & \textbf{Owner} & \textbf{Deliverable} \\
  \midrule
  Wk 6 & K-Fold cross-validation ($k=10$) + temporal holdout evaluation + AUC-ROC curves & R1+R2 & Validation chapter + graphs \\
  Wk 6 & End-to-end integration test: live scan $\rightarrow$ detection $\rightarrow$ alert $\rightarrow$ block $\rightarrow$ MTD rotation & All & Demo recording \\
  Wk 7 & Final report: CRISP-DM structure, literature review, ethical considerations & All & Final report PDF \\
  Wk 7 & French abstract (r\'{e}sum\'{e}), English abstract, architecture diagram, slides & R5 & Submission package \\
  \bottomrule
\end{tabular}
\end{table}


% ============================================================
%  13. TECHNICAL STACK SUMMARY
% ============================================================
\section{Technical Stack Summary}

\begin{table}[h!]
\centering
\small
\begin{tabular}{L{4cm} L{10cm}}
  \toprule
  \textbf{Component} & \textbf{Tool / Library} \\
  \midrule
  OS / Virtualization    & Ubuntu Server (IDS + honeypot), Kali Linux (attacker), VirtualBox / GNS3 \\
  Programming Language   & Python 3.10+ \\
  Traffic Capture        & Scapy with BPF filters, sliding window aggregation \\
  ML Libraries           & Scikit-learn, XGBoost, imbalanced-learn, Joblib, Pandas, NumPy \\
  Web Dashboard          & Flask + Socket.IO + D3.js \\
  Blocking Engine        & Linux \texttt{iptables} DROP/REDIRECT via Python subprocess, TTL-based auto-unblock \\
  MTD Engine             & Custom Python service: dynamic port reassignment + config push to legitimate nodes \\
  Honeypot Framework     & Cowrie (SSH/Telnet) or Dionaea (multi-service) via Docker \\
  Attack Simulation      & Nmap (SYN/ACK/UDP/XMAS/Connect), Masscan, custom Scapy scripts \\
  Version Control        & Git + GitHub --- feature branches per module \\
  \bottomrule
\end{tabular}
\end{table}


% ============================================================
%  14. EXPECTED OUTCOMES & SUCCESS CRITERIA
% ============================================================
\section{Expected Outcomes \& Success Criteria}

\begin{enumerate}[nosep,leftmargin=1.5em]
  \item Model achieves F1-Score $\geq 88\%$ on CICIDS2017 PortScan test subset.
  \item An aggressive Nmap SYN stealth scan is detected and an alert is triggered
        within 8 seconds of initiation.
  \item A slow-rate scan (1 probe/second) is correctly identified within the
        90-second analysis window.
  \item The attacker profile correctly identifies scan type (SYN / UDP / ACK /
        XMAS) and includes honeypot contact history where applicable.
  \item At least 60\% of scanning source IPs interact with at least one honeypot
        node during a full scan campaign.
  \item The MTD engine completes a port rotation cycle within 30 seconds of a
        confirmed detection trigger.
  \item Following MTD rotation, continued probing of vacated ports by the same
        source IP is correctly classified as post-scan activity with no false
        negative.
  \item The dashboard reflects scanning events, honeypot interactions, and MTD
        rotation status with a maximum update delay of 5 seconds.
  \item Automated blocking is enforced within 15 seconds of a confirmed detection.
  \item Legitimate traffic generates a false positive rate below 6\% across all
        test scenarios, including during active MTD rotation cycles.
\end{enumerate}


% ============================================================
%  15. CRISP-DM COVERAGE CHECKLIST
% ============================================================
\section{CRISP-DM Report Chapter Mapping}

\begin{table}[h!]
\centering
\small
\begin{tabular}{C{0.8cm} L{5cm} L{8cm}}
  \toprule
  & \textbf{CRISP-DM Phase} & \textbf{Covered By} \\
  \midrule
  $\square$ & Business Understanding        & Sections 1--3 of this BRD \\
  $\square$ & Data Understanding            & Section 5 (datasets) + Section 6 (features) \\
  $\square$ & Data Preparation              & SMOTE, time-window aggregation (Section 7) \\
  $\square$ & Modeling                       & RF + XGBoost + Isolation Forest (Section 7) \\
  $\square$ & Evaluation                    & K-Fold, temporal holdout, comparison table (Section 9) \\
  $\square$ & Deployment                    & Architecture (Sec 10), Honeypot (Sec 8), MTD (Sec 8) \\
  \midrule
  $\square$ & Literature Review             & \textit{Existing IDS approaches, ML for network security} \\
  $\square$ & Ethical Considerations         & \textit{Privacy, automated blocking implications} \\
  $\square$ & French Abstract (R\'{e}sum\'{e}) & \textit{Required for final submission} \\
  $\square$ & English Abstract              & \textit{Required for final submission} \\
  \bottomrule
\end{tabular}
\end{table}


% ============================================================
%  16. FUTURE WORK / PERSPECTIVES
% ============================================================
\section{Perspectives / Future Work}

The following items are \textbf{out of scope} for the current project but are
documented as promising research directions for future iterations:

\begin{itemize}[nosep,leftmargin=1.5em]
  \item \textbf{AI-Integrated MTD:} Feed the MTD state table as a live input
        to the ML model, adding an ``MTD Port Delta'' feature that captures the
        difference between the probed port and the current assigned port. This
        requires generating MTD-aware training data through controlled
        simulation --- infeasible within the current 7-week timeline.
  \item \textbf{Deep Learning models:} Evaluate LSTM or 1D-CNN architectures
        for temporal pattern recognition across multi-window sequences.
  \item \textbf{Distributed IDS deployment:} Extend the system from a single
        LAN to a multi-subnet architecture with federated model updates.
  \item \textbf{High-interaction honeypot:} Deploy a full sandboxed OS
        (instead of Cowrie/Dionaea) to capture more advanced attacker behavior,
        including malware samples and lateral movement attempts.
\end{itemize}


\end{document}
